% Transcription — fonds Alexandre Grothendieck, shelfmark 161-5, batch 1 (pages 1-20).
% Established from the facsimile archives/batches/161-5/batch-01.pdf (a mirror of
% https://grothendieck.umontpellier.fr/161-5.pdf), per the skill
% transcribe-grothendieck. The batch file carries no cover sheet: PDF page N is
% archivists' page N, as the pencilled numbers bottom-left confirm.
%
% Pass: Opus 5.5 (claude-opus-5-5), 2026-09-22 — first pass, unchecked against the pages by a human.
% Revised 2026-09-23 (Opus 5.5): header only — « Four of the five point to later
% work » replaced by an exact count (all five name Sen or Springer; four name a
% publication). Body unchanged.
%
% What the batch is. Nothing in it is manuscript in the full sense: it is two
% printed offprints by J.-P. Serre, annotated in red ink, and a folder cover.
%
%   1      The folder (a manila chemise). Its tab carries « Barsotti-Tate » in
%          pencil, written up the edge; nothing else but the archivists' « (5) ».
%          No \page; the word heads the file and a \note says where it comes
%          from and that the hand is not surely his.
%   2-16   Offprint: J.-P. Serre, « Sur les groupes de Galois attachés aux
%          groupes p-divisibles », Proceedings of a Conference on Local Fields
%          (NUFFIC Summer School, Driebergen 1966), Springer 1967, pp. 118-131.
%          Printed, not transcribed. Pages 9, 15 and 16 carry annotations and
%          appear for those alone, each anchored by a \note naming the printed
%          statement; pages 2-8 and 10-14 carry nothing in any hand and are
%          absent. For the modernised reading: the printed footnote 1 (page 3)
%          says that the idea of replacing p-adic Lie algebras by the
%          corresponding algebraic groups « m'avait d'ailleurs été suggérée par
%          Grothendieck il y a plusieurs années, en liaison avec sa théorie des
%          motifs » — printed, Serre's words, so not in the transcription.
%   17-20  Offprint: extract from the Annuaire du Collège de France, 68e année
%          (1968-1969), Résumé des cours de 1967-1968 — Serre, « Algèbre et
%          géométrie », pp. 47-49 on pages 18-20 (page 17 is the extract's
%          title page). Printed, not transcribed; pages 19 and 20 are annotated.
%          The résumé goes on past page 20 (the « Séminaire » section begins at
%          the foot of p. 49, and a list of publications shows through from the
%          verso), so batch 2 opens in the middle of it.
%
% The annotations. Five, all in the same red ink and the same hand; one blue
% correction of a printed sign (page 15) is noted but not attributed. All five
% name S. Sen or T. Springer: four name Sen (page 9 by name alone, « Annals »
% on pages 15 and 19, « Inventions » on page 16), the fifth (page 20) Springer's
% book on Jordan algebras « récemment paru chez Springer-Verlag »; four of the
% five thus name a publication, page 9 none.
% Those references are transcribed exactly as written and not completed here;
% for whoever maintains the modernised reading, they are presumably Sen's
% Invent. Math. 17 (1972) and Ann. of Math. 97 (1973) papers and Springer's
% « Jordan algebras and algebraic groups » (1973), which would put the
% annotations after 1972, whereas the inventory gives 1967-1969 for the
% offprints. That is inference, not something the pages state; it stays out of
% \dating. Whether the red hand is Grothendieck's was not settled: nothing in
% it is characteristic enough to decide, and the file says so once, at the top.
%
% Notation: the anchoring notes follow Serre's printed notation (G_alg, I_alg,
% G_V, H_V^0 for the neutral component, V_C(n)); H_V^0, printed with a small
% raised circle in the Annuaire, is set H_V^\circ.
\documentclass[11pt,a4paper]{article}
% Le préambule commun à toutes les transcriptions du fonds Grothendieck.
%
% Il tient en une idée : **l'incertitude se compose, elle ne se résout pas.**
% Une transcription qui lisse un mot douteux a détruit la seule chose pour
% laquelle on la faisait — un lecteur doit pouvoir savoir, à chaque endroit,
% ce qui était sur le feuillet et ce qui a été ajouté. Les six macros qui
% suivent sont donc l'appareil critique, et non de la décoration.
%
% Les mêmes commandes sont reconnues par `scripts/render.mjs`, qui produit la
% vue de lecture du site. Une macro ajoutée ici doit l'être aussi là-bas,
% faute de quoi le rendu échouera bruyamment — ce qui est voulu : un rendu
% silencieusement faux est pire qu'un rendu absent.

\ProvidesPackage{grothendieck}[2026/08/08 Transcriptions du fonds Grothendieck]

\usepackage{amsmath,amssymb,amsthm}
\usepackage{mathrsfs}
\usepackage{stmaryrd}
% Les diagrammes commutatifs sont partout dans ce fonds : c'est la langue
% même de la géométrie algébrique de Grothendieck, et une transcription qui
% les rendrait en prose ne transcrirait rien.
\usepackage{tikz-cd}
% Français par défaut — c'est la langue des feuillets — mais l'anglais est
% chargé aussi : la traduction et le résumé sont des documents anglais, et la
% typographie française (espaces avant les deux-points) y serait une faute.
% Ces documents basculent par \selectlanguage{english} en tête de corps.
\usepackage[english,french]{babel}
\usepackage{fontspec}
\usepackage{geometry}
\geometry{a4paper,margin=25mm,marginparwidth=28mm}
\usepackage{marginnote}
\usepackage{xcolor}
% ulem, not soul. `soul`'s \st and \ul are built on a character-by-character
% scan that XeLaTeX's font machinery breaks: the document fails with an
% undefined control sequence rather than degrading. ulem does the same two jobs
% and is Unicode-engine safe. `normalem` stops it hijacking \emph, which the
% transcriptions use for Grothendieck's own emphasis.
\usepackage[normalem]{ulem}
\usepackage{hyperref}
\hypersetup{colorlinks=true,linkcolor=black,urlcolor=black,pdfborder={0 0 0}}

% --- Métadonnées du lot -----------------------------------------------------
% Reprises telles quelles par le script de rendu, qui les affiche en tête de
% la vue de lecture. Elles fixent ce que le fichier prétend couvrir : sans
% elles, une transcription flotte sans point d'attache dans un fonds de seize
% mille feuillets.

\newcommand{\folder}[1]{\gdef\grothFolder{#1}}
\newcommand{\batch}[1]{\gdef\grothBatch{#1}}
\newcommand{\leaves}[2]{\gdef\grothFirst{#1}\gdef\grothLast{#2}}
\newcommand{\foldertitle}[1]{\gdef\grothFoldertitle{#1}}
\gdef\grothFolder{?}\gdef\grothBatch{?}\gdef\grothFirst{?}\gdef\grothLast{?}\gdef\grothFoldertitle{}

% --- Le résumé --------------------------------------------------------------
%
% Il ouvre la lecture modernisée, sous le titre « Résumé », et il est composé
% en petit. Ce n'est pas un ornement : c'est le seul passage du document qui
% ne soit pas des mathématiques, et un lecteur doit pouvoir le reconnaître
% avant de l'avoir lu — puis le sauter s'il connaît déjà le sujet, ou s'y
% arrêter s'il ne le connaît pas. Le filet qui le suit marque la reprise du
% corps.

\newenvironment{resume}
  {\small\setlength{\parskip}{0.45em}}
  {\par\medskip\hrule\medskip}

% --- Mots-clés ---------------------------------------------------------------
%
% En anglais, en clôture du résumé de la lecture modernisée : le vocabulaire
% moderne sous lequel chercher ce que les feuillets construisent. Le manifeste
% du site (`npm run manifest`) les extrait pour étiqueter la cote entière —
% cette ligne est la source unique des tags, il n'y a pas de fichier à côté.
\newcommand{\keywords}[1]{\par\smallskip\noindent
  {\footnotesize\textsc{Keywords} --- \textit{#1}}\par}

% --- L'appareil critique ----------------------------------------------------

% Le numéro de feuillet, en marge. C'est l'ancre qui permet de régler un
% désaccord sur une formule en regardant *un* feuillet plutôt qu'une chemise
% de six cents. Le site s'en sert aussi pour tourner les pages du fac-similé
% quand on fait défiler la transcription.
\newcommand{\leaf}[1]{%
  \marginnote{\footnotesize\color{gray!70}\textsf{#1}}%
  \label{leaf:#1}\ignorespaces}

% Illisible : on ne devine pas. Un mot inventé qui se lit comme les autres est
% le pire résultat possible.
\newcommand{\ill}{\textcolor{red!60!black}{[\,\ldots\,]}}

% Lecture douteuse : le mot est donné, et signalé comme incertain.
\newcommand{\uncertain}[1]{\uline{#1}}

% Ajout de l'éditeur — une lettre manquante, un mot sous-entendu.
\newcommand{\add}[1]{\textcolor{blue!50!black}{[#1]}}

% Biffé par Grothendieck lui-même. Ce qu'il a rayé fait partie du document :
% une rature dit souvent où la pensée a changé de direction.
\newcommand{\struck}[1]{\sout{#1}}

% Note du transcripteur, sur l'état matériel du feuillet.
\newcommand{\note}[1]{\par\smallskip{\small\color{gray!60!black}\textit{[#1]}}\par\smallskip}

% Note marginale de Grothendieck, distincte du corps du texte.
\newcommand{\marginal}[1]{\par\smallskip\noindent\hspace{1em}%
  {\small\color{gray!60!black}#1}\par\smallskip}

% --- Titre ------------------------------------------------------------------

\AtBeginDocument{%
  \begin{center}
    {\large\bfseries Fonds Alexandre Grothendieck --- cote n\textsuperscript{o}~\grothFolder}\\[2pt]
    {\itshape\grothFoldertitle}\\[4pt]
    {\small feuillets \grothFirst--\grothLast\ (lot \grothBatch)}\\[2pt]
    {\footnotesize\color{gray!60!black}Transcription établie d'après le fac-similé numérisé
     par l'Université de Montpellier.\\
     Les lectures incertaines sont soulignées, les illisibles notées [\,\ldots\,],
     les ajouts entre crochets.}
  \end{center}
  \vspace{1em}}

\endinput


\folder{161-5}
\batch{1}
\pages{1}{20}
\dating{1967-1969 — le groupe « Dossiers rassemblés par Grothendieck sur différentes thématiques » (161-1 à 162-6) est daté [après 1961-vers 1977]}
\watermark{Édition de démonstration}
\foldertitle{Barsotti-Tate : tirés à part annotés (1967-1969), notes manuscrites (s.d.).}

\begin{document}

\section*{Barsotti-Tate}

\note{Ce mot est écrit au crayon, en long, sur l'onglet de la chemise (page 1),
qui ne porte rien d'autre ; la main n'est pas sûrement la sienne. Les pages 2 à
20 sont deux tirés à part imprimés de J.-P. Serre. Le texte imprimé n'est pas de
sa main et n'est pas transcrit ; ne sont relevées que les annotations
manuscrites, avec l'indication de l'énoncé qu'elles visent. Elles sont toutes à
l'encre rouge et d'une même main, qu'on n'a pas pu identifier avec certitude
comme la sienne.}

\subsection*{Tiré à part : J.-P. Serre, « Sur les groupes de Galois attachés aux groupes $p$-divisibles » (1967)}

\page{9}
\note{Tiré à part, imprimé : J.-P. Serre, « Sur les groupes de Galois attachés
aux groupes $p$-divisibles », \emph{Proceedings of a Conference on Local
Fields}, NUFFIC Summer School, Driebergen 1966, Springer-Verlag, 1967,
p. 118-131 (pages 2 à 16 du dossier). Seules les pages 9, 15 et 16 portent des
annotations.}

\marginal{En fait, on a $M = I_{\mathrm{alg}}$ (Sen)}
\note{en marge de gauche, en regard de la Remarque 3 qui clôt le § 3 (p. 124 de
l'article). $K_{nr}$ y est l'extension non ramifiée maximale de $K$ contenue
dans $\bar{K}$, $I$ l'image de $\mathrm{Gal}(\bar{K}/K_{nr})$ dans $G$ (le
sous-groupe d'inertie), $V$ est un $I$-module semi-simple, et l'imprimé énonce
$M \subset I_{\mathrm{alg}}$, où $M$ est le groupe de Mumford-Tate du module
galoisien $V$ et $I_{\mathrm{alg}}$ l'enveloppe algébrique de $I$ ; il le tire
du théorème 2 appliqué sur le corps de base $\widehat{K}_{nr}$. L'annotation
remplace l'inclusion par une égalité.}

\page{15}
\note{Dans la démonstration du lemme 3 (p. 130 de l'article), l'inégalité
imprimée « pour $v(x) < c_2 = c_1/(p^h - 1)$ » porte un signe $<$ repassé à
l'encre bleue sur une retouche blanche ; ce qu'il recouvre ne se lit pas, et
rien n'indique de quelle main est la correction.}

\marginal{voir Sen (\uncertain{Annals})}
\note{en marge de gauche, en face d'une parenthèse tracée le long de la
Remarque 1 qui suit le Théorème 5. Le Théorème 5 : pour $F$ un groupe formel
sur $A$, de dimension 1 et de hauteur finie $h$, avec
$\mathrm{End}(F) = \mathbb{Z}_p$, l'image $G$ de $\mathrm{Gal}_K$ dans
$\mathrm{Aut}(V)$ est ouverte. La Remarque 1 : sans l'hypothèse
$\mathrm{End}(F) = \mathbb{Z}_p$, on peut encore déterminer l'algèbre de Lie
$\mathfrak{g}$ de $G$, qui est le commutant de $\mathrm{End}(F)$ dans $V$, par
une démonstration analogue mais sensiblement plus compliquée. Le dernier mot
est écrit vite et abrégé ; la lecture « Annals » est celle de la page 19, où il
se lit en clair.}

\page{16}
\marginal{démontré par S. Sen (Inventions)}
\note{en marge de droite, au bout d'un crochet qui embrasse la question finale
de la Remarque 3 (p. 131 de l'article). Le groupe d'inertie $I$ y est muni de
deux filtrations,
\[
v(g) \geqslant n \iff (g-1)(T) \subset p^n T, \qquad
w(g) \geqslant n \iff g \in I^n
\]
($I^n$ le $n$-ième groupe de ramification en numérotation supérieure), et
l'imprimé demande si, $e_K$ désignant l'indice de ramification absolu de $K$,
on a $w = e_K \cdot v + O(1)$ ; le crochet couvre aussi la phrase suivante, qui
pose la même question pour tous les groupes de Lie $p$-adiques qui sont des
groupes de Galois.}

\subsection*{Tiré à part : J.-P. Serre, « Algèbre et géométrie », résumé des cours de 1967-1968 (Annuaire du Collège de France)}

\page{19}
\note{Extrait imprimé de l'\emph{Annuaire du Collège de France}, 68e
année (1968-1969), Résumé des cours de 1967-1968 : J.-P. Serre, « Algèbre et
géométrie », à partir de la p. 47 (pages 17 à 20 du dossier ; la suite passe au
lot suivant). $V$ y est une représentation $p$-adique du groupe de Galois $G$
de $\bar{K}/K$, $G_V$ son image, $H_V$ l'enveloppe algébrique de $G_V$ et
$h : \mathbb{G}_m/C \to \mathrm{GL}_V/C$ l'homomorphisme défini par la
décomposition $V_C = \sum V_C(n)$. Seules les pages 19 et 20 sont annotées.}

\marginal{ceci a été démontré par S. Sen (Annals)}
\note{en marge de gauche, en face d'un crochet qui embrasse un alinéa du point
2) (p. 48 de l'Annuaire). L'alinéa précédent caractérise, lorsque $G_V$ est
résoluble, $H_V^{\circ}$ comme le plus petit sous-groupe algébrique $R$ de
$\mathrm{GL}_V$ sur $\mathbb{Q}_p$ tel que $R/C$ contienne
$h(\mathbb{G}_m/C)$, l'analogue $p$-adique du groupe de Mumford-Tate.
L'alinéa marqué conjecture que cette caractérisation vaut sans supposer $G_V$
résoluble, et observe qu'on pourrait le prouver si l'on savait que
$V_C = V_C(0)$ entraîne que $G_V$ est fini.}

\page{20}
\marginal{voir à ce sujet le livre de T. Springer sur les alg. de Jordan
récemment paru chez Springer-Verlag}
\note{en marge de droite, au bout d'un crochet en face de la dernière phrase du
point 5) (p. 49 de l'Annuaire). Le point 5) suppose que $V$ ne fait intervenir
que les exposants 0 et 1, $V_C = V_C(0) + V_C(1)$, avec
$n_0 = \dim_C V_C(0)$ et $n_1 = \dim_C V_C(1)$ (cas du module attaché à un
groupe $p$-divisible, $n_1$ étant sa dimension) ; pour $n_1 = 1$ et $V$ simple,
$G_V$ contient un sous-groupe ouvert de $\mathrm{Aut}_E(V)$, $E$ le commutant
de $H_V^{\circ}$. La phrase marquée dit que le cours s'est achevé par un essai
de classification dans le cas $n_1 > 1$ ($V$ absolument simple sur
$H_V^{\circ}$), fondé sur ce que $H_V^{\circ}/C$ contient un tore de dimension
1 n'ayant que deux poids.}

\end{document}
